\section{Discussion}\label{sec:discussion}

The results separate three questions that steering studies often conflate: whether an intervention preserves coherent behavior, whether it changes the target behavior, and only then where that change is implemented.

\paragraph{The instrument must be checked before the mechanism is probed.}
Our naive sweep produced a result that would superficially appear publishable, a benchmark gain together with a $\rho \approx 1$ depth verdict, from a vector that multiplies held-out perplexity \NaivePplRatio-fold. Nothing in the standard evaluation pipeline flags this. Teacher-forced multiple choice is structurally insensitive to generation collapse, and nothing in a ratio statistic prevents dividing one artifact by another. The two preconditions we impose, a denominator bounded away from zero at a coherence-gated operating point and an aligned-versus-random contrast, are inexpensive to compute, and our results indicate that they should constitute standard practice for causal claims built on steering interventions. On the present evidence, steering-strength sweeps that report only benchmark deltas are not interpretable.

\paragraph{Where steering is vocabulary-level, it does not work.}
Once that gate is applied, the CAA honesty vector presents the first substantive lesson. It moves honesty-coded vocabulary strongly, over a logit of curated honest-shift at the answer position, while leaving truthfulness flat to slightly negative on held-out questions. An evaluation that scored word choice, persona, or self-reported honesty would classify this outcome as a success, whereas TruthfulQA with a held-out split does not. The two families were built from different supervision: contrast prompts that instruct the model to \emph{behave} honestly, versus statements labeled by what is \emph{true}. The outcome is consistent with each construction tracking what its supervision specified, a surface register in one case and something upstream of the readout in the other, though what the mass-mean direction actually represents is not identified by our experiments. For safety applications the operational distinction is central: an intervention that makes a model sound more honest without making it more truthful is a failure mode that word-choice- or persona-based evaluation can misread as alignment progress.

\paragraph{Where the steering effect works, it is predominantly not vocabulary suppression.}
The mass-mean construction supplies the complementary result. For this construction with a real behavioral effect, the deflationary hypothesis that motivated this work is at most a minority of the story. On Llama-3 the mass-mean vector's truthfulness gain survives certified excision of its direct readout onto every aligned token set we tested, under norm matching, across three token-set constructions, under paraphrase shift, and in free-form generation. In-domain, the aligned subspace is not distinguished from arbitrary subspaces: its excision is statistically indistinguishable from random excisions of equal rank, and the contrast interval rules out the large aligned-specific cost that suppression predicts. The effect is carried predominantly by the processing that downstream layers perform on the injected direction rather than by what the tested unembedding rows read from it directly. Section~\ref{sec:multimodel} reports the same pattern on Mistral and Llama-2, where the majority ($\rho = \MistralMmRho$ and $\LlamaTwoMmRho$ respectively) is again downstream, although aligned excision removes approximately one quarter of the Mistral effect. The conclusion that survives across the three models with a real effect is comparative, not absolute: wherever a steering vector improves truthfulness, most of that improvement survives aligned direct-readout excision.

\paragraph{Implications for readout-based interpretability.}
This sequence of results changes how the vectors themselves should be read. The lens trajectories show that the aligned-64 parallel component produces negligible word-level shift, while the certified-orthogonal component reproduces the full trajectory; the per-prompt word-level depth ratios confirm that most of each vector's vocabulary effect is produced downstream even when the aligned direct path is available. This is a concrete, quantified instance of the non-identifiability that \citet{venkatesh2026nonid} establish in general: the component of a steering vector visible to the unembedding is a poor guide to what the vector does. Methods that interpret steering vectors by reading them through the vocabulary projection, such as logit-lens decoding of steering directions or vocabulary-space vector arithmetic, isolate only a component that is neither necessary nor sufficient for the full Llama-3 effect.
